\chapter{怎样从硬件事实推导内核对象}

上一章已经把硬件链路单独讲清楚：CPU 只能从当前寄存器和程序计数器继续执行；地址要经过页表、TLB、cache 和内存属性；设备通过 MMIO、队列、DMA、中断和 completion 与 CPU 协作。现在往上走一层：怎样从这些硬件事实推导内核对象。

task、地址空间、VMA、fd、request、wait queue 这些名字不是先验规定出来的。它们解决的都是同一个问题：硬件只给瞬时状态和低层事件，OS 必须把它们组织成可恢复、可检查、可等待、可取消、可释放的长期状态。只要这条线不清楚，后面讲调度、虚拟内存、文件系统和驱动时，就会像在背结构体字段。

先从一个非常普通的动作开始。一个进程往日志文件追加一行：

\begin{lstlisting}
int fd = open("app.log", O_APPEND | O_WRONLY);
write(fd, user_buf, len);
\end{lstlisting}

直觉上，这像是“把 \code{user\_buf} 里的字节交给磁盘”。但硬件并不认识“进程”“文件”“字符串”“追加写”这些概念。硬件只认识当前指令、寄存器里的数、虚拟地址、页表权限、cache line、设备寄存器、DMA 地址、中断向量和 completion 队列。于是内核必须把这次写日志拆成一串可验证的状态转换：当前程序怎样停下来，用户地址怎样验证，fd 怎样变成文件对象，文件偏移怎样确定，数据怎样进入 page cache 或块请求，设备完成怎样回来，失败后哪些状态已经提交，哪些状态必须回滚。

\begin{keyidea}
内核对象的第一原则是：把硬件给出的低层状态，变成可恢复、可检查、可等待、可释放的状态。进程不是“正在运行的代码”四个字，而是 CPU 现场、地址空间、文件表、信号、等待原因和调度状态的组合；I/O 请求也不是“调用设备函数”，而是 buffer、设备地址、request id、completion 和等待者的组合。
\end{keyidea}

把这条路径先压成一张总表：

\begin{longtable}{p{0.18\linewidth}p{0.36\linewidth}p{0.34\linewidth}}
\toprule
硬件事实 & 如果只写成普通函数会怎样 & 由此建立的内核对象 \\
\midrule
CPU 只能从当前现场继续执行 & 进入内核后覆盖寄存器，返回用户态时找不到原位置 & trap frame、task、内核栈 \\
用户地址只是虚拟地址 & 用户传坏指针会让内核读错页、越权读，或跨页时只检查了一半 & address space、VMA、PTE、用户拷贝边界 \\
\code{fd} 只是一个小整数 & 另一个线程 close 后对象可能被释放，权限和文件偏移也无处保存 & fd table、open file description、引用计数 \\
设备只暴露寄存器和队列 & 把设备当函数调用会忙等、丢完成、无法取消和超时 & driver、request、descriptor ring \\
完成事件异步回来 & 中断只说“有事件”，不能说明哪个请求完成、结果是什么 & completion、request id、wait queue \\
DMA 会让设备直接访问内存 & CPU 和设备同时改 buffer，或设备写到已经释放的页 & buffer ownership、DMA mapping、IOMMU 约束 \\
\bottomrule
\end{longtable}

这张表也是后面章节的阅读索引。调度章讲 task 为什么能被切走再回来；内存章讲 address space、VMA、PTE 为什么分层；系统调用章讲用户指针为什么不能直接用；设备章讲 request、descriptor 和 completion 怎样接到硬件；文件章讲 fd 怎样引用文件对象；安全章讲谁有权触碰这些对象。读到任何 OS 名词时，都应该问一句：它是在补哪条硬件事实缺失的语义？

\section{一、从一次写日志看硬件边界}

先写一个错误但很自然的版本：

\begin{lstlisting}
write(fd, user_buf, len):
  disk.write(user_buf, len)
  return len
\end{lstlisting}

这个版本的问题不是“代码太短”，而是它把五个硬件边界都抹掉了。

第一，CPU 正在执行用户程序，不会凭空同时运行内核代码。用户态进入内核时，硬件只提供一个受控入口。入口必须保存返回地址、栈指针、状态寄存器和参数寄存器。否则内核处理完以后，不知道回到哪条用户指令，也不知道用户原来传了什么参数。

第二，\code{user\_buf} 不是物理内存地址，也不是可信内核指针。它只是当前进程地址空间里的一个虚拟地址。它可能跨页，可能后半段不可读，可能指向尚未调入的 \code{mmap} 文件页，也可能在另一个线程里被并发 \code{munmap}。内核不能因为起始地址看起来正常，就把 \code{len} 字节当成一整段可靠内存。

第三，\code{fd} 不是文件本身。它只是当前进程文件表里的一个小整数索引。内核要用它找到 open file description，也就是“这个打开实例”的状态：访问权限、当前文件偏移、append 标志、引用计数、底层 inode 或设备对象。没有这个对象，两个线程同时写同一个 fd 时，谁更新偏移、谁持有引用、close 以后对象何时释放，都说不清。

第四，磁盘不是函数。高速存储设备通常靠队列工作：CPU 在内存里放描述符，写 MMIO doorbell 通知设备；设备经 DMA 读写内存；完成后写 completion 并触发中断或等待轮询。设备返回不是沿 C 调用栈返回，而是稍后通过硬件事件回来。

第五，失败不是一种。用户地址错误、文件不可写、磁盘超时、设备 reset、只写入一部分、数据只到 page cache、数据到设备缓存但未持久化，这些都不同。内核对象必须记录“推进到哪一步、哪些字节已经进入哪个层次、哪个请求还在飞、谁在等它、失败证据在哪里”。

现在把例子走细一点。假设用户执行：

\begin{lstlisting}
write(fd, user_buf, 6000);
\end{lstlisting}

\code{user\_buf} 从用户页 A 的偏移 3000 开始，长度 6000 字节，因此会跨过页 A、页 B，并进入页 C。页 A 已经映射且可读；页 B 当前 PTE 不 present，但 VMA 说明它是合法的文件映射，可以缺页调入；页 C 不在任何 VMA 内。fd 指向一个普通日志文件，打开时带 \code{O\_APPEND}。

朴素版本会说：“检查起点合法，然后拷 6000 字节。”这会错。正确路径必须按边界推进：

\begin{longtable}{p{0.14\linewidth}p{0.42\linewidth}p{0.32\linewidth}}
\toprule
阶段 & 内核和硬件发生什么 & 为什么要形成对象 \\
\midrule
入口 & syscall 保存用户 RIP/RSP、状态位和参数寄存器，切到内核入口和内核栈 & trap frame 让内核能返回、重启或把错误发给正确任务 \\
fd 查找 & 当前 task 的 fd table 把小整数 \code{fd} 解析成 open file description，并临时增加引用 & 防止并发 close 释放对象，也保存写权限和 append 语义 \\
权限和偏移 & 检查文件是否可写、是否 append、是否受 mount 或安全策略限制；append 写要在提交点取文件尾 & 文件对象保存“这次打开”的权限和偏移状态 \\
用户拷贝页 A & CPU 按用户页表读取页 A 尾部，复制到内核 buffer 或 page cache 页 & 用户地址访问仍可能 fault，不能当普通内核 memcpy \\
用户拷贝页 B & PTE 不 present，CPU 进入 page fault；内核查 VMA 后把页调入，再重试拷贝 & VMA 是“这个 fault 是否合法”的软件策略，PTE 是硬件证据 \\
遇到页 C & 地址不在任何 VMA 内，内核不能继续读取 & 请求对象要记录已经接收多少字节，返回短写或错误语义 \\
进入文件层 & 已接收的字节进入 page cache，或构造成块层请求；相关页标 dirty & page cache 把用户 API 返回和设备持久化分开 \\
提交设备 & 若需要回写，块层分配 request，驱动填 descriptor、DMA mapping、request id，写 doorbell & 设备只认队列、DMA 地址、长度和控制位 \\
睡眠等待 & 当前 task 从 running 变成 blocked，挂到 request 的 wait queue & CPU 不应忙等慢设备；等待原因必须可恢复 \\
完成返回 & 设备写 completion，驱动按 tag 找回 request，记录结果并唤醒等待 task & 中断只是入口，request 才说明哪次请求完成 \\
\bottomrule
\end{longtable}

这条路径里，每个对象都有必要性。trap frame 解决“回哪里”；task 解决“谁正在请求，谁要睡眠”；address space 和 VMA 解决“用户地址是否合法”；file object 解决“fd 对应哪个打开实例”；page cache 解决“文件内容在内存里的版本”；request 解决“异步设备请求的身份”；wait queue 解决“谁在等这个条件”；completion 解决“设备完成事实怎样回到具体请求”。

再看页 C。内核已经成功接收页 A 尾部和页 B 的一部分之后，发现页 C 无效。此时不能简单说“整个 write 失败”，也不能简单说“6000 字节都成功”。普通文件、管道、socket、直接 I/O 对短写和错误的语义不同，但共同要求是：内核必须知道哪些字节已经进入了哪个内核对象，哪些还没有越过边界。这就是为什么后面文件章会反复讲 short write，为什么系统调用章会讲 \code{copy\_from\_user} 不是普通 \code{memcpy}。

再看 append。没有 \code{O\_APPEND} 时，写入使用 open file description 里的当前文件偏移；写完后偏移前进。带 \code{O\_APPEND} 时，写入位置应当在提交时取文件尾。若两个线程同时写同一个日志文件，内核必须把“取文件尾、写入、更新大小或偏移”放在文件系统定义的同步边界里，否则两条日志可能互相覆盖。硬件只会执行 load/store，OS 对象必须补出“这是一个打开实例，这是一次追加写，这是已经提交的文件大小变化”。

同一次 \code{write} 放到不同路径里，硬件边界还会继续分叉。很多初学者把“写文件”直接想成“CPU 把字节送到磁盘”，这会把 page cache、直接 I/O、设备 DMA 和崩溃一致性混成一件事。先看四种常见路径：

\begin{longtable}{p{0.18\linewidth}p{0.36\linewidth}p{0.34\linewidth}}
\toprule
路径 & 表面看起来 & 真正跨过的硬件边界 \\
\midrule
page cache 命中 & 字节很快被复制到内核页，\code{write} 返回 & CPU 从用户页读，写入 page cache 页；设备请求可能稍后才出现 \\
page cache 未命中或部分块写 & 先准备文件页，再把用户字节合进去 & 可能需要从设备读旧块，填补用户没有覆盖的部分，再标 dirty \\
直接 I/O & 绕过 page cache，让设备直接读用户数据 & 用户页要 pin，建立 DMA/IOMMU 映射，completion 前不能释放或改作他用 \\
append 与崩溃边界 & 多个线程都说“追加一行” & 文件大小、目录项、日志事务、设备缓存 flush 的提交点必须分清 \\
\bottomrule
\end{longtable}

第一条路径是 page cache 命中。文件对应的缓存页已经在内存中，内核把用户字节复制到这些页，给页打上 dirty 标记，然后 \code{write} 可以较快返回。这里的完成只说明“内核已经接收这些字节，并把文件的内存版本改了”。它不说明磁盘已经持久化。硬件层面实际发生的是 CPU load 用户页、CPU store 内核页、cache line 被改脏；块设备可能还没有看到任何 request。于是 page cache 页必须保存 dirty 状态、文件偏移、页索引和引用计数，否则后台回写线程将来不知道哪些页要写到哪个文件位置。

第二条路径是 page cache 未命中，尤其是部分块写。假设文件系统块大小是 4096 字节，用户只覆盖某个块中间 100 字节。如果这个块原内容不在 page cache，内核不能凭空把一个新页的 100 字节写出去，因为同一块其余 3996 字节仍然要保留旧内容。常见做法是先把旧块读进 page cache，或通过文件系统规则确认其余部分可以视为零，再把用户字节合进去。这里一个 \code{write} 可能先产生读 request，再产生未来的写 request。对象模型必须能表达“这页正在读入”“这页已经 uptodate”“这页被用户写脏”“这页等待回写”。如果只写成 \code{disk.write(user\_buf, len)}，就无法解释为什么一次写有时会先读磁盘。

第三条路径是直接 I/O。用户用 \code{O\_DIRECT} 或数据库用自己的缓存时，内核可能让设备 DMA 直接从用户页读数据。此时用户地址检查还不够，内核还要 pin 住参与 I/O 的物理页，防止它们在设备访问期间被换出、释放或重新映射给别人。然后驱动通过 DMA API 把这些页变成设备可见的 DMA 地址，必要时写入 IOMMU 页表。completion 到达前，request 保存的不是一份普通指针，而是一组仍被设备引用的页、一组 DMA 映射和一个方向标记：这是设备读内存还是写内存。若进程在直接 I/O 期间退出，内核也不能立即释放这些页，必须等 request 完成或被设备确认取消。

第四条路径是 append、并发和崩溃。两个线程同时执行：

\begin{lstlisting}
write(fd, "A\n", 2);
write(fd, "B\n", 2);
\end{lstlisting}

如果二者都先读到同一个文件尾，再各自写入，就可能互相覆盖。文件对象和文件系统锁必须把“取尾部位置并分配写入范围”做成一个受控动作。再往下，系统突然断电时，用户看到 \code{write} 返回并不等价于数据已经越过设备缓存落到介质上；\code{fsync} 返回也依赖文件系统日志、块层 flush、设备对 flush/FUA 的实现。这里需要的对象不仅是 open file 和 page cache，还有 journal transaction、dirty inode、writeback request 和错误回报位置。后面文件系统章会继续展开，但此处先记住一点：硬件只提供写 cache line、发 request、收到 completion 这些事实，OS 必须把它们组合成“追加”“可见”“持久化”这些接口语义。

把四条路径放在一起，就能看出本章的主线：同一个 API 名字下面有多层完成边界。page cache 命中路径主要考验用户拷贝、页状态和延迟回写；部分块写路径暴露读改写；直接 I/O 路径暴露 DMA 和 IOMMU；append 路径暴露并发提交和崩溃恢复。内核对象不是为了把代码写复杂，而是为了让这些不同硬件事实不会互相冒充。

\begin{longtable}{p{0.18\linewidth}p{0.38\linewidth}p{0.32\linewidth}}
\toprule
如果没有这个对象 & 直接后果 & 对应的真实问题 \\
\midrule
没有 trap frame & 系统调用或中断后无法恢复用户指令流 & 任务不能被安全暂停和恢复 \\
没有 task & 寄存器现场、地址空间、文件表和等待状态无法归属 & 调度器不知道切谁、唤醒谁 \\
没有 VMA & PTE 不 present 时不知道是懒分配还是非法地址 & 缺页处理只能盲目失败或盲目放行 \\
没有 open file description & fd 无法保存权限、偏移和引用关系 & 并发 write/close 语义崩掉 \\
没有 request & 设备完成回来时找不到原始 buffer 和等待者 & 中断只能说明“有事”，不能完成 API \\
没有 wait queue & 线程要么忙等，要么睡了以后没人知道它等什么 & 慢设备会浪费 CPU 或丢唤醒 \\
\bottomrule
\end{longtable}

所以本章不是给结构体起名，而是把一次普通写入拆成对象形成过程。对象不是为了“面向对象设计”，而是为了保存跨越硬件边界的事实：谁拥有、谁等待、谁能访问、谁已经提交、谁还可能失败。

\section{二、为什么要从 CPU 入口建立 task 和 trap frame}

先只看 CPU，不看文件和设备。用户程序执行到这里：

\begin{lstlisting}
write(fd, user_buf, len);
next_line();
\end{lstlisting}

进入内核前，CPU 的软件可见状态大致包括：

\begin{lstlisting}
RIP    = address of syscall instruction or next return point
RSP    = top of user stack
RAX    = syscall number
RDI    = fd
RSI    = user_buf
RDX    = len
RFLAGS = condition flags, interrupt-related state, mode bits
\end{lstlisting}

这些名字不用一次背完。关键是：CPU 当前只知道一组寄存器和下一条指令地址。用户程序执行 \code{syscall} 后，硬件会进入内核配置好的入口，并保存返回用户态所需的一部分状态。内核入口代码随后把更多通用寄存器保存成一份结构化现场。这份现场通常叫 trap frame。

trap frame 第一次出现时，可以把它理解成“进入内核那一刻的证据包”。它至少回答四个问题：

\begin{longtable}{p{0.2\linewidth}p{0.36\linewidth}p{0.32\linewidth}}
\toprule
问题 & trap frame 里的证据 & 后续用途 \\
\midrule
从哪里来 & 用户 RIP、用户代码段或特权级状态 & 返回用户态、信号处理、调试、错误报告 \\
怎么回去 & 用户 RSP、返回状态位、中断相关状态 & 恢复用户栈和受控状态 \\
带了什么参数 & 系统调用号和参数寄存器 & 派发到具体系统调用并验证参数 \\
失败在哪里 & faulting RIP、fault address、error code & 缺页、非法指令、权限错误和重启逻辑 \\
\bottomrule
\end{longtable}

如果不保存现场，内核很快会把用户寄存器覆盖掉：

\begin{lstlisting}
bad_syscall_entry:
  use rdi, rsi, rdx for kernel work
  call filesystem_write
  return_to_user   // old fd, user_buf, len, and return point are gone
\end{lstlisting}

这不是抽象风险。内核自己的 C 函数也需要寄存器传参，也会使用栈，也会调用其他函数。若入口不先保存用户现场，内核一旦开始工作，原来的用户状态就会被正常的内核执行覆盖。trap frame 把“用户现场”从 CPU 瞬时状态变成内存里的对象，后面的检查、调度、信号和返回才有基础。

但 trap frame 还不够。系统调用期间可能发生定时器中断，当前线程可能睡眠，调度器可能选择另一个线程运行。此时需要一个更大的归属对象：task。task 可以先理解成“一个可被调度和恢复的执行流”。它不是单独的寄存器集合，而是把一组长期状态绑在一起：

\begin{lstlisting}
Task:
  saved CPU context
  kernel stack
  state: running, runnable, blocked, stopped
  address space
  file table
  signal state
  wait reason and wait queue link
  accounting and scheduling metadata
\end{lstlisting}

用一次系统调用睡眠来走完整路径。用户线程调用 \code{write}，进入内核，内核发现设备请求需要等待，于是把当前 task 标成 blocked，并挂到 request 的 wait queue 上。调度器选择另一个 runnable task，保存旧 task 的内核执行上下文，恢复新 task 的上下文。等设备完成后，中断处理函数把旧 task 从 wait queue 移回 runqueue。未来某个时间片里，调度器恢复这个 task，它继续从系统调用内部的睡眠点之后执行，最终返回用户态。

\begin{longtable}{p{0.15\linewidth}p{0.43\linewidth}p{0.3\linewidth}}
\toprule
时刻 & 状态变化 & 为什么需要 task \\
\midrule
用户执行 & CPU 当前寄存器属于用户线程 & 这些状态要归属到某个可调度实体 \\
系统调用入口 & trap frame 记录用户现场，内核栈开始工作 & 进入内核后仍要知道将来回哪个用户线程 \\
提交 I/O & request 记录 buffer、长度、tag 和等待者 & 等待关系要能从 request 回到 task \\
睡眠 & task 状态从 running 变成 blocked，离开 CPU & CPU 可以运行其他 task，原 task 不丢失 \\
中断完成 & completion 找到 request，再找到等待 task & 硬件事件变成对具体 task 的唤醒 \\
再次调度 & task 被放回 runnable 并恢复内核上下文 & 系统调用从睡眠点继续，而不是从头再来 \\
返回用户态 & trap frame 恢复用户寄存器和返回位置 & 用户程序看到一个普通 \code{write} 返回值 \\
\bottomrule
\end{longtable}

这张表解释了为什么“进程/线程”不能只说成“一个程序”。程序是代码和数据；task 是当前执行流的可恢复状态。两个线程可以运行同一份代码、共享同一个地址空间和文件表，但有不同的寄存器现场、栈、等待原因和调度状态。OS 需要的是后者，因为 CPU 只能在具体现场之间切换。

再加入异步中断。假设用户线程正在循环计算，没有调用系统调用：

\begin{lstlisting}
while (true) {
  work();
}
\end{lstlisting}

若没有定时器中断，这个线程可以一直占着 CPU。定时器到期后，硬件在指令边界进入内核入口。入口保存被打断现场，调度器给当前 task 记账，必要时把它放回 runqueue，再选择另一个 task。这里没有用户主动配合，task 仍然成立，因为硬件中断入口和 trap frame 让内核能把“当前 CPU 正在执行的现场”抓成对象。

\begin{keyidea}
task 是 CPU 硬件事实的内核化：CPU 只能执行一个当前现场，OS 就把每个可恢复现场做成 task；CPU 入口只能给一次瞬时 trap，OS 就把入口证据保存成 trap frame；CPU 会被中断打断，OS 就把等待原因和调度状态也放进 task。
\end{keyidea}

再把入口类型分清。进入内核并不只有系统调用一种。系统调用是用户程序主动请求服务；page fault 是 CPU 在执行 load/store/fetch 时发现地址翻译或权限不满足；timer interrupt 是时钟设备要求内核获得控制权；device interrupt 是外设报告队列或状态有变化。它们都进入内核，但原因、保存的证据和后续动作不同。

\begin{longtable}{p{0.18\linewidth}p{0.35\linewidth}p{0.35\linewidth}}
\toprule
入口 & 触发原因 & trap frame 之外还要找什么对象 \\
\midrule
system call & 用户代码执行受控入口指令，请求内核服务 & 当前 task、fd table、address space、系统调用参数 \\
page fault & MMU 无法完成某次取指、读或写 & fault address、error code、VMA、PTE、可能的物理页 \\
timer interrupt & 定时器到期，CPU 被异步打断 & 当前 task 的运行时间、runqueue、调度策略 \\
device interrupt & 设备报告队列、completion 或错误状态 & driver、device queue、request tag、wait queue \\
\bottomrule
\end{longtable}

系统调用入口的特点是“主动而同步”。用户代码知道自己正在调用 \code{write}，参数按 ABI 放在寄存器里，返回值会回到同一条执行流。内核主要做两件事：把不可信参数变成内核对象，例如 fd 变成 open file；把可能阻塞的动作变成 task 的等待状态，例如等待 I/O completion。

page fault 的特点是“当前指令没有完成”。例如内核在 \code{copy\_from\_user} 中读 \code{user\_buf}，CPU 发现页 B 的 PTE 不 present，于是进入 page fault。trap frame 记录的是被 fault 打断的位置，硬件还提供 fault address 和错误码。内核查 VMA 后若能补页，就安装 PTE 并返回到原指令附近重试；若不能补救，就把错误转成系统调用失败，或给用户进程发信号。这里 trap frame 不只是返回用的，它还让内核知道“这次 fault 属于哪个 task、发生在用户态还是内核态、能不能恢复”。

timer interrupt 的特点是“用户或内核代码没有主动请求”。CPU 正在执行某个 task，定时器突然让控制权转入内核。内核保存现场后，可以决定当前 task 的时间片已经用完，把它从 running 放回 runnable，再从 runqueue 里挑另一个 task。若没有 trap frame，异步打断会破坏被打断代码；若没有 task，调度器只知道“CPU 上有寄存器”，不知道这些寄存器属于谁、已经运行多久、该放回哪个队列。

device interrupt 的特点是“完成事实在设备侧”。设备中断到来时，当前 CPU 上运行的 task 未必是当初提交 I/O 的 task。甚至提交 I/O 的 task 可能在另一个 CPU 上睡眠。中断处理函数不能靠当前 task 来猜结果，必须读取设备 completion，按 tag 找 request，再通过 request 找等待者或回调。这一层关系非常重要：中断入口保存的是“被打断的当前现场”，completion 保存的是“设备完成了哪次请求”。二者不是同一个对象。

还要解释内核栈。用户态有用户栈，为什么进入内核后不能继续用用户栈？因为用户栈位于用户地址空间，权限和映射都由用户进程影响。用户可以把栈指针设到非法地址，也可以让栈页在进入内核后触发缺页，还可以把栈内容设计成让内核覆盖关键数据。如果内核直接在用户栈上保存返回地址、局部变量和锁状态，就等于把内核控制流交给了不可信内存。所以每个 task 需要受内核保护的内核栈。系统调用、中断、缺页处理进入内核后，内核代码在这块栈上运行；task 睡眠时，这条内核调用链也随 task 一起保存。

\begin{longtable}{p{0.2\linewidth}p{0.35\linewidth}p{0.33\linewidth}}
\toprule
如果继续用用户栈 & 可能发生什么 & 内核栈解决什么 \\
\midrule
用户 RSP 指向非法页 & 入口保存现场时再次 fault，甚至无法可靠报告错误 & 入口有可信保存位置 \\
用户 RSP 指向共享映射 & 另一个线程或进程可修改内核临时数据 & 内核调用链不受用户写入影响 \\
用户栈空间不足 & 深层文件系统或驱动调用覆盖用户数据或触发嵌套错误 & 每个 task 有受控大小和保护页 \\
系统调用中睡眠 & 用户可能修改栈上“内核状态” & 睡眠中的内核帧保存在内核地址空间 \\
\bottomrule
\end{longtable}

CPU 入口还会发生在内核态。内核正在执行文件系统代码时，设备中断可能到来；内核正在做用户拷贝时，page fault 可能到来；内核拿着某个锁时，定时器中断可能到来。于是 trap frame 还要记录被打断时的特权级和状态。若中断来自用户态，返回路径可能经过调度、信号投递和用户态恢复；若中断来自内核态，返回路径通常要回到被打断的内核代码，并且必须遵守“哪些锁已持有、哪些栈帧仍有效、哪些操作不可重入”。这就是为什么 OS 入口代码往往比普通函数调用复杂：它不仅要进来，还要能在各种嵌套和异步条件下正确出去。

\section{三、为什么要从地址翻译建立 address space、VMA 和 PTE}

现在看 \code{user\_buf}。用户传给内核的是一个虚拟地址。虚拟地址是进程自己看到的地址编号，例如 \code{0x7fff00001234}。它不直接说明真实内存条上的位置，也不说明当前是否允许读写。硬件里的 MMU 查页表，把虚拟页翻译成物理页，并检查权限。页表项叫 PTE，也就是 page table entry。

可以先把 PTE 看成硬件能读懂的一条记录：

\begin{lstlisting}
PTE:
  present: is this virtual page currently mapped?
  physical page number
  readable / writable / executable
  user-accessible or kernel-only
  dirty / accessed and other status bits
\end{lstlisting}

但是 PTE 只回答“当前硬件能不能直接完成这次访问”。它不回答“如果不能完成，内核该不该补救”。这个补救策略来自 VMA。VMA 可以理解成 address space 里的一个合法区间说明：这段虚拟地址从哪里来，允许读写执行吗，背后是匿名内存、文件映射、栈，还是设备映射。

\begin{longtable}{p{0.18\linewidth}p{0.4\linewidth}p{0.3\linewidth}}
\toprule
层次 & 它保存什么 & 谁使用它 \\
\midrule
address space & 一个进程看到的虚拟地址世界，包含页表根和 VMA 集合 & task 切换、缺页处理、用户地址检查 \\
VMA & 一段虚拟区间的策略：范围、权限、来源、增长规则 & 内核缺页处理和 \code{mmap}/\code{munmap} \\
PTE & 硬件当前可执行的映射：present、物理页、权限位 & MMU page walk 和 TLB 填充 \\
TLB & 最近 PTE 翻译结果的硬件缓存 & CPU 快速完成 load/store \\
\bottomrule
\end{longtable}

用前面的跨页 \code{user\_buf} 走一遍：

\begin{lstlisting}
page A: PTE present, user readable
page B: PTE not present, VMA says file mapping and readable
page C: no VMA covers this address
\end{lstlisting}

拷贝页 A 时，硬件翻译通过，CPU 从用户页读出字节。拷贝页 B 时，硬件发现 PTE 不 present，于是进入 page fault。注意，PTE 不 present 不等于一定非法。内核要查 VMA：若 VMA 说明这是一段合法文件映射，内核可以分配物理页、从文件读入数据、安装 PTE，然后重试原来的用户读取。拷贝页 C 时，没有任何 VMA 覆盖这个地址，内核就不能补页，因为这不是“暂时没准备好”，而是“这段地址不属于进程”。

\begin{longtable}{p{0.18\linewidth}p{0.38\linewidth}p{0.32\linewidth}}
\toprule
访问结果 & 硬件看到什么 & 内核怎样解释 \\
\midrule
页 A 成功 & PTE present 且用户可读，TLB 或 page walk 给出物理页 & 直接复制这一页内剩余字节 \\
页 B 缺页 & PTE not present，CPU 记录 fault address 和错误码 & 查 VMA，合法则补页后重试 \\
页 C 非法 & 没有 VMA，或者 VMA 不允许读 & 停止拷贝，返回错误或短写语义 \\
只读页写入 & PTE present 但 writable 位不允许 & 可能是 COW，也可能是非法写 \\
NX 执行失败 & PTE 不允许执行 & 安全隔离，不应改成可执行后放行 \\
\bottomrule
\end{longtable}

这就是 address space、VMA、PTE 三层同时存在的原因。只有 PTE，不知道 not-present 是懒分配、文件页、换出页，还是非法地址。只有 VMA，硬件无法每条 load/store 都进内核询问策略。只有 TLB，修改页表后旧翻译还可能继续被 CPU 使用。OS 必须维护这三层的一致关系。

把几种 fault 分开看，会更清楚。初学者常把 page fault 理解成“程序访问了非法内存”。这只对一部分 fault 成立。page fault 更准确的说法是：CPU 按当前页表和权限位无法完成这次内存访问，于是把证据交给内核。内核再根据 VMA 和对象状态判断，这是正常的按需工作，还是应该终止进程。

\begin{longtable}{p{0.19\linewidth}p{0.36\linewidth}p{0.33\linewidth}}
\toprule
情况 & 硬件报告 & 内核解释 \\
\midrule
匿名内存第一次写 & PTE not present，地址落在可写匿名 VMA 内 & 分配一页清零物理页，安装 PTE，然后重试 \\
文件映射第一次读 & PTE not present，地址落在可读 file-backed VMA 内 & 从文件对应位置读页，填入 page cache 或映射页，再重试 \\
换出页被访问 & PTE not present，但软件记录说明页在 swap 中 & 从 swap 读回，恢复 PTE，再重试 \\
COW 写入 & PTE present 但不可写，VMA 允许写且页被共享 & 复制物理页，给当前进程安装私有可写 PTE \\
真正越权写 & PTE 或 VMA 不允许写 & 返回 \code{EFAULT} 或向用户进程发送信号 \\
NX 取指 & PTE 不允许执行，CPU 正在取指 & 阻止把数据页当代码执行，通常发送信号 \\
\bottomrule
\end{longtable}

匿名内存第一次写是最温和的 fault。用户调用 \code{malloc} 后，库可能只向内核申请一段虚拟地址范围，真实物理页直到第一次写入才分配。CPU 写这个地址时发现 PTE not present，进入内核。内核查到地址属于可写匿名 VMA，就分配一页物理内存、清零、安装 PTE。返回后，原来的写指令重新执行，用户程序完全可以感觉不到这次 fault。这里 VMA 保存的是“这段地址应该存在，只是物理页还没准备”；PTE 保存的是“现在硬件还不能直接访问”。

文件映射第一次读类似，但来源不同。用户用 \code{mmap} 把文件映射到地址空间，读取某个尚未在内存中的页时，内核要根据 VMA 找到底层文件、页偏移和权限，把文件内容读入页，再建立映射。若此时文件被截短，访问落在新文件大小之外，VMA 仍然可能覆盖这段虚拟地址，但底层文件对象已经没有对应内容。很多系统会把这种错误报告成类似 \code{SIGBUS} 的总线错误。这个例子说明：VMA 合法不代表底层对象永远能提供数据；VMA 还必须和文件大小、页缓存状态、文件系统错误一起解释。

COW，也就是 copy-on-write，能说明 PTE 权限位不只是安全检查，也是一种延迟复制机制。父进程 \code{fork} 后，父子进程先共享同一批物理页，PTE 都标成只读。某一方写入时，CPU 因为 writable 位不允许而 fault。内核查 VMA 后发现这不是非法写，而是 COW 写，于是分配新页、复制旧内容、给当前进程安装可写 PTE。另一个进程仍然指向旧页。若只看硬件错误码，会以为这是“写只读页”；只有结合 VMA 和页引用计数，内核才知道这是“该复制了”。

NX fault 则是安全边界。NX 可以理解成 PTE 上的“不可执行”权限。栈、堆、普通数据页通常不应当被当作代码取指。CPU 从这些页取指时进入 fault，内核不能因为“页存在、可读”就放行执行。这里 PTE 把读写执行拆成不同权限，VMA 记录这段区域的预期用途，安全策略再决定是否允许把某段内存改成可执行。后面讲进程安全时，这条边界会直接连接到 W\^{}X、JIT、动态链接器和代码注入防护。

TLB 会让问题再多一层。CPU 为了不每次访存都走完整页表，会把最近的虚拟页到物理页翻译缓存到 TLB。内核修改 PTE 后，如果旧翻译还留在某个 CPU 的 TLB 中，那个 CPU 可能继续按旧权限访问旧物理页。例如内核把某页从可写改成只读，用来启动 COW；如果没有让相关 CPU 刷掉旧 TLB 项，另一个 CPU 可能继续写共享页，COW 就被绕过。再比如内核撤销某段用户映射后，旧 TLB 项若继续有效，用户代码可能访问已经释放并被重新分配的物理页。于是 address space 对象不只保存页表根，还要知道哪些 CPU 可能缓存了这个地址空间的翻译，什么时候需要本地或远程 TLB shootdown。

\begin{longtable}{p{0.18\linewidth}p{0.38\linewidth}p{0.32\linewidth}}
\toprule
PTE 变化 & 如果 TLB 没同步 & 后果 \\
\midrule
present 变 not present & CPU 继续用旧翻译访问物理页 & 访问已释放页，破坏新对象 \\
writable 变只读 & CPU 继续写旧可写映射 & COW 或写保护失效 \\
user 变 kernel-only & 用户态继续访问旧用户映射 & 权限隔离失效 \\
物理页号改变 & CPU 继续访问旧物理页 & 读写错对象，数据难以追踪 \\
\bottomrule
\end{longtable}

用户拷贝和普通用户访存还有一个差别：fault 发生在内核代码执行期间。内核执行 \code{copy\_from\_user} 时，当前特权级在内核态，但被访问地址属于用户地址空间。若这个访问 fault，不能简单按“内核 bug”崩掉，也不能像用户态 fault 那样随便重启任意指令。内核必须知道这段代码是一个允许 fault 的用户拷贝窗口，fault 可恢复时继续，不可恢复时把错误返回给系统调用。很多内核会给这类拷贝点配套异常表：某条内核指令访问用户地址失败时，跳到预先登记的修复位置。对象层面看，这仍然是 trap frame、address space、VMA 和系统调用返回路径一起工作。

再看一个具体反例。假设内核为了省事，在系统调用入口只检查 \code{user\_buf} 起点：

\begin{lstlisting}
bad_copy_from_user(user_buf, len):
  if address_is_valid(user_buf):
    memcpy(kernel_buf, user_buf, len)
\end{lstlisting}

这段代码在跨页时会错。起点所在页合法，不代表后续每一页合法。更严重的是，页表可能在拷贝过程中发生变化：另一个线程可以 \code{munmap} 某段地址，文件映射页可能第一次访问才调入，栈页可能按需增长。可靠的用户拷贝必须以“可能 fault、可能部分成功、必须能恢复”为前提设计。

\begin{lstlisting}
copy_from_user_like(dst, user_addr, len):
  copied = 0
  while copied < len:
    check current user page
    if page fault can be resolved:
      resolve and retry
    if access is illegal:
      stop and report copied/error
    copy bytes until page boundary
    copied += bytes_this_page
\end{lstlisting}

这里的循环不是实现细节，而是语义边界。它告诉后面的系统调用章节：用户指针永远不能先验可信；告诉内存章节：缺页是硬件事件和 VMA 策略的交汇；告诉文件章节：短写和错误返回来自“请求推进到一半时遇到不可继续边界”。

地址空间还解释了为什么 task 切换不只是换寄存器。若调度器从进程 A 切到进程 B，只恢复寄存器还不够；同一个虚拟地址在两个进程里可以指向不同物理页。内核还要切换页表根，或者切换带地址空间标识的翻译上下文。否则进程 B 的 \code{user\_buf} 可能被按进程 A 的页表解释。

\begin{longtable}{p{0.18\linewidth}p{0.38\linewidth}p{0.32\linewidth}}
\toprule
切换内容 & 少做会怎样 & 对象来源 \\
\midrule
通用寄存器和 PC & 新 task 从错误位置继续，参数和局部值损坏 & CPU 现场 \\
内核栈 & 两个 task 的内核调用链互相覆盖 & task 私有内核执行状态 \\
页表根或地址空间标识 & 用户虚拟地址按错误进程解释 & address space \\
TLB 相关状态 & CPU 继续使用旧地址空间的翻译 & PTE/TLB 缓存一致性 \\
文件表引用 & fd 指向的对象归属混乱或被提前释放 & task 或进程资源表 \\
\bottomrule
\end{longtable}

地址空间还会和 DMA 发生关系。设备 DMA 使用的是设备可见地址，不一定等于 CPU 物理地址，更不等于用户虚拟地址。直接 I/O 中，用户传入的 \code{user\_buf} 必须先被分解成一批用户虚拟页；内核检查 VMA 和 PTE，确保这些页允许这次方向的 I/O；再把对应物理页 pin 住；最后通过 IOMMU 建立设备可见映射。IOMMU 可以理解成“给设备用的地址翻译和权限检查”。没有它，设备一旦拿到错误 DMA 地址，可能读写任意物理内存；有了它，驱动仍然要正确建立、同步和撤销映射。

\begin{longtable}{p{0.18\linewidth}p{0.38\linewidth}p{0.32\linewidth}}
\toprule
地址名字 & 谁使用 & 不能混淆的原因 \\
\midrule
用户虚拟地址 & 用户程序和内核用户拷贝代码 & 只在某个 address space 中有意义 \\
内核虚拟地址 & 内核代码 & 可能映射同一物理页，也可能没有直接映射用户页 \\
物理页号 & CPU 页表和内核页管理器 & 设备未必直接用这个编号访问 \\
DMA 地址 & 设备和 IOMMU & 生命周期受 DMA mapping 控制 \\
\bottomrule
\end{longtable}

如果一个进程提交直接 I/O 后马上 \code{munmap(user\_buf)}，用户虚拟地址可以从 address space 中消失，但被 pin 的物理页不能立刻释放，因为设备 request 仍然引用它们。等 completion 到达、驱动同步 DMA、撤销 IOMMU 映射并放掉页引用后，这些页才真正回到普通内存管理。这个例子把第三节和第四节连起来：地址空间解决 CPU 访问的合法性，request 和 DMA mapping 解决设备访问期间的所有权。

从硬件角度看，地址空间是页表和权限检查；从内核角度看，地址空间是一个进程能触碰哪些虚拟地址、这些地址来自哪里、出现 fault 时怎样补救的对象。二者必须同时成立。

\section{四、为什么要从设备请求建立 request、wait queue 和 completion}

假设内核已经把用户数据安全地接收到内核对象里，下一步要让设备写出去。错误版本会写成：

\begin{lstlisting}
device.write(kernel_buf, len);
\end{lstlisting}

真实设备通常暴露三类东西：MMIO 寄存器、内存中的队列、完成事件。MMIO 寄存器是 CPU 通过特殊内存属性访问的设备状态槽；队列是 CPU 和设备都能读写的一段内存；完成事件可能是 completion queue 中的一项，也可能伴随中断。

先看最小串口模型：

\begin{lstlisting}
STATUS bit0: tx_ready
DATA:        write one byte to start sending
CTRL bit0:   enable transmitter
\end{lstlisting}

朴素轮询版是：

\begin{lstlisting}
while ((read_mmio(STATUS) & tx_ready) == 0) {
  // keep spinning
}
write_mmio(DATA, byte);
\end{lstlisting}

这个模型能说明 MMIO，但不适合通用 OS 的慢设备路径。若设备 10 ms 后才 ready，CPU 这 10 ms 都在原地读状态寄存器，不能运行别的 task。于是内核需要把“现在不能继续，但将来条件满足后要回来”表示成等待关系。等待关系不能只存在栈上，因为当前 task 会离开 CPU；它必须挂到内核对象上。

wait queue 第一次出现时，可以把它理解成“等同一个条件的 task 链表”。等待串口 ready 的 task 可以挂在串口对象的 wait queue 上；等待块请求完成的 task 可以挂在 request 的 wait queue 上；等待锁的 task 可以挂在锁或 futex 对应的队列上。睡眠不是消失，而是 task 状态从 running 变成 blocked，并记录它等的条件。

高吞吐设备会进一步使用 request。request 是一次异步设备工作的内核身份牌。它至少保存：

\begin{lstlisting}
Request:
  id or tag
  target device queue
  operation: read or write
  buffer list and length
  DMA mapping
  current state: new, submitted, completed, failed, canceled
  error code
  waiter or completion callback
\end{lstlisting}

用 NVMe 写请求走一遍：

\begin{longtable}{p{0.15\linewidth}p{0.43\linewidth}p{0.3\linewidth}}
\toprule
阶段 & 发生什么 & request 保存什么 \\
\midrule
创建 & 文件系统或块层把脏页转成块 I/O 请求 & 目标 LBA、长度、方向、关联页 \\
映射 & 驱动通过 DMA API 建立设备可见地址 & DMA 地址、方向、长度和待撤销映射 \\
填描述符 & CPU 在 submission queue 写 opcode、DMA 地址、tag & tag 把设备完成对回 request \\
敲门铃 & CPU 写 MMIO doorbell，通知设备队列 tail 前进 & request 状态变成 submitted \\
睡眠或异步返回 & 同步调用把 task 挂到 wait queue；异步路径记录回调或完成对象 & 谁在等、完成后通知谁 \\
设备 DMA & 设备读取描述符并搬运数据 & buffer 不能释放或复用 \\
写 completion & 设备把 tag、status、phase 写到 completion queue & request 结果等待驱动收割 \\
中断处理 & 驱动读取 completion，按 tag 找到 request & 状态变成 completed 或 failed \\
唤醒与回收 & 唤醒等待 task，撤销 DMA mapping，释放引用 & 生命周期结束条件清楚 \\
\bottomrule
\end{longtable}

这里最容易混的是中断和 completion。中断只是 CPU 入口，意思是“有设备事件需要处理”。completion 才是设备写下的结果，告诉驱动哪个 tag 完成、状态码是什么。若只根据中断返回用户态，内核不知道完成的是哪次 request；若只看 completion 不处理等待者，任务会永远睡着。

再看 DMA。DMA 不是“更快的 memcpy”，而是设备获得了直接访问内存的能力。于是 buffer 所有权必须明确：

\begin{longtable}{p{0.18\linewidth}p{0.4\linewidth}p{0.3\linewidth}}
\toprule
阶段 & 谁拥有 buffer & 规则 \\
\midrule
提交前 & CPU/内核拥有 & 可以填数据、建描述符、建立 IOMMU 映射 \\
提交后完成前 & 设备拥有或共享拥有 & CPU 不能释放、复用或无序改写；必须遵守 cache/DMA 同步规则 \\
completion 后 & 结果回到驱动，但还没完全释放 & 先读取状态、同步 cache、撤销 DMA mapping \\
回收后 & CPU/内核重新拥有 & 页或 buffer 才能回到分配器或上层对象 \\
\bottomrule
\end{longtable}

如果驱动在 completion 前释放 buffer，设备迟到 DMA 可能写坏新对象。如果驱动没有建立正确 IOMMU 映射，设备可能无法访问 buffer，或者越界访问不该碰的页。如果驱动没有在 descriptor 写完后再敲 doorbell，设备可能读到旧描述符。request 对象把这些边界放到一处：提交前准备，提交后禁止释放，完成后按顺序回收。

失败路径更能说明 request 的必要性。设备可能超时。超时不代表设备一定没做；它只代表“在规定时间内没有观察到 completion”。此时内核不能马上释放 buffer，因为设备可能稍后还会 DMA。可靠恢复通常要停止队列、屏蔽中断或 quiesce 设备、尝试 abort request、必要时 reset 控制器、确认没有 in-flight DMA 后再撤销映射。没有 request 列表，内核不知道哪些请求还在飞；没有 request state，内核不知道哪些能重试、哪些已经失败、哪些不能再释放。

为了避免把设备统一想成“磁盘”，再看三个设备路径。它们的规模不同，但都要求同一组 OS 对象。

\begin{longtable}{p{0.16\linewidth}p{0.39\linewidth}p{0.33\linewidth}}
\toprule
设备路径 & 典型动作 & 所需对象 \\
\midrule
PIO 串口发送一个字节 & CPU 读状态寄存器，ready 后写 DATA 寄存器 & wait queue、设备锁、寄存器访问顺序 \\
网卡接收一个包 & 设备 DMA 到预先给出的 buffer，写 completion 或更新 ring & buffer pool、descriptor ring、packet object、NAPI/poll 状态 \\
NVMe 写多个块 & CPU 填 submission queue，设备 DMA 用户或 page cache 页，写 completion queue & request tag、DMA mapping、queue depth、timeout/reset 状态 \\
\bottomrule
\end{longtable}

PIO，也就是 programmed I/O，是 CPU 自己一字节或一字地读写设备寄存器。它简单，但暴露了等待问题。串口发送缓冲满时，CPU 可以忙等，也可以让当前 task 睡到“发送缓冲有空位”这个条件上。若使用中断，设备在变 ready 时触发中断，驱动唤醒 wait queue。这里没有复杂 DMA，仍然需要 wait queue，因为硬件 ready 信号和用户 \code{write} 的返回不是同一时刻发生的。

网卡接收路径能说明“设备先写内存，CPU 后处理”。驱动会提前准备一批接收 buffer，把它们的 DMA 地址放进 receive ring。网卡收到包后，把包内容 DMA 到某个 buffer，再在 completion 或 descriptor 状态位里写明长度、校验状态和 buffer 编号。中断到来后，驱动并不是从设备寄存器里读出整个包，而是沿着 ring 找到哪些 descriptor 已完成，把 buffer 变成内核网络包对象，再补充新的空 buffer 给设备。若 buffer pool 没有对象，网卡可能丢包；若驱动过早复用 buffer，上层协议可能读到被新包覆盖的数据。

NVMe 这类块设备则把队列深度推高。一个队列里可以同时有几十、几百甚至更多 request 在飞。CPU 不能用“我刚提交的就是下一个完成的”这种假设，因为设备可能重排、合并或不同请求耗时不同。request tag 的作用就在这里：提交时把 tag 写入 descriptor，completion 回来时带回 tag，驱动据此找到原 request。没有 tag 或等价身份，异步高并发设备就无法正确匹配结果。

再看“走 cache”和“不走 cache”的差别。这里的 cache 不只指 CPU cache，也包括 page cache 和设备内部缓存。普通 buffered write 先进入 page cache，后续由回写线程合并、排序并提交块请求；直接 I/O 绕过 page cache，但仍然经过 CPU cache 一致性、IOMMU 和设备队列；设备内部可能还有易失缓存，completion 只说明设备接受或完成了命令，不一定说明掉电后仍可恢复。不同缓存层解决的问题不同：

\begin{longtable}{p{0.19\linewidth}p{0.36\linewidth}p{0.33\linewidth}}
\toprule
缓存层 & 它让什么变快 & OS 必须额外记录什么 \\
\midrule
CPU cache & CPU 读写内存更快 & DMA 前后的可见性、内存屏障、cache 同步 \\
page cache & 文件读写复用内存页，合并小 I/O & dirty、uptodate、writeback、文件偏移到页索引的关系 \\
块层队列缓存/合并 & 多个请求排序、合并、限流 & request 状态、队列深度、公平性和超时 \\
设备内部缓存 & 设备接收写入后更快返回 & flush/FUA、持久化错误回报、断电边界 \\
\bottomrule
\end{longtable}

假设日志写入走 buffered path：用户字节先复制到 page cache，\code{write} 返回；几秒后回写线程把脏页变成 block request；设备 completion 回来后，页的 writeback 状态清掉。这个路径下用户态很快继续执行，但崩溃时最近数据可能还在内存或设备缓存中。若应用调用 \code{fsync}，内核要把相关脏页、元数据和日志事务推到设备承诺的持久化边界，并处理设备可能晚报的写错误。

假设同一日志写入走 direct I/O：用户页被 pin，request 保存 scatter-gather 列表，设备 DMA 直接读取这些页。这里 page cache 不承担数据缓存，但 CPU cache 仍然存在。如果平台是 cache coherent DMA，硬件能保证 CPU cache 和设备 DMA 看到的一致性；如果平台是 non-coherent DMA，驱动必须在设备读内存前把 CPU 写过的数据清到内存，在设备写内存后让 CPU cache 失效或同步。否则设备可能读到旧数据，CPU 也可能读不到设备刚写的新数据。

内存屏障也来自这个边界。CPU 和编译器可能为了性能重排普通内存写；设备则按它观察到的内存和 MMIO 顺序行动。驱动填 descriptor 时，必须保证 descriptor 内容对设备可见之后，才写 doorbell。否则设备看到 doorbell 后去读 descriptor，可能读到半旧半新的内容。抽象成代码就是：

\begin{lstlisting}
fill descriptor fields
make descriptor writes visible to device
write MMIO doorbell
\end{lstlisting}

中间那句在不同架构上可能是内存屏障、DMA 同步或有序 MMIO 规则。它不是“玄学优化”，而是在告诉硬件：先看到队列内容，再看到通知。没有这层顺序，request 对象里保存的状态和设备实际执行的命令可能对不上。

再看丢中断。设备 completion 已经写到内存，但中断因为屏蔽、合并或硬件问题没有及时送到 CPU。若驱动只依赖“中断来了才检查 completion”，request 可能一直挂起。许多高性能驱动会把中断和轮询结合：中断负责提示“可能有活”，poll 循环负责批量收割 completion；超时定时器负责在长时间没有完成时检查 in-flight request。这样即使单个中断丢失，completion ring 中的结果仍可能被后续轮询发现。这里 wait queue 等的是 request 状态，而不是“某一次中断一定发生”。

设备 reset 是最复杂的失败路径之一。reset 不是把结构体清零这么简单。reset 前可能已有 request submitted，设备可能已经 DMA 了一半，completion 可能在 reset 过程中丢失，队列内存可能还留着旧 phase 位。可靠恢复至少要回答：

\begin{longtable}{p{0.25\linewidth}p{0.55\linewidth}}
\toprule
问题 & 必须有的状态 \\
\midrule
哪些请求已经提交给设备 & in-flight request list 和 queue generation \\
哪些 buffer 仍可能被设备访问 & DMA mapping、pinned pages、设备 quiesce 结果 \\
哪些请求可以重试 & request 操作类型、幂等性、文件系统上层语义 \\
哪些请求必须失败上报 & error code、等待 task、回调对象 \\
旧 completion 怎样区分 & tag、phase、queue generation 或 reset epoch \\
\bottomrule
\end{longtable}

这张表说明“复杂硬件”不是因为名词多，而是因为异步边界多。CPU 可以继续运行，设备可以稍后 DMA，中断可能合并，completion 可能乱序，reset 可能跨过半完成状态。OS 要把这些边界收束成对象生命周期，否则用户看到的 \code{read}/\code{write}/\code{fsync} 就没有可靠语义。

\begin{longtable}{p{0.18\linewidth}p{0.38\linewidth}p{0.32\linewidth}}
\toprule
事件 & 不能怎么处理 & 正确依赖的对象 \\
\midrule
completion 正常到达 & 不能只说“中断来了”就返回成功 & completion tag、request status、wait queue \\
设备超时 & 不能立刻释放 DMA buffer & in-flight request list、timeout state、reset protocol \\
用户取消 I/O & 不能假设设备已经停止访问内存 & request cancel state、设备 abort 或完成确认 \\
进程退出 & 不能释放仍被设备持有的用户页 & pinned page list、request refcount、DMA mapping \\
设备 reset & 不能沿用 reset 前的队列状态 & queue generation、request recovery、completion cleanup \\
\bottomrule
\end{longtable}

这就是设备章和文件章会反复出现 request、completion、wait queue、refcount、timeout 的原因。它们不是工程噪声，而是异步硬件进入 OS 语义后的最低对象集合。

\section{五、第一阶段映射：从硬件事实到内核对象模型}

现在把前四节收束成一个最小内核模型。这个模型不是 Linux 完整结构体图，而是后文所有章节反复使用的检查框架：每个内核对象都必须回答“保存什么事实、谁能引用它、谁在等待它、它何时提交、它何时释放”。

\begin{lstlisting}
Machine
  CPUs[]:
    current task
    trap/interrupt entry state
    current address-space root

  PhysicalPages[]:
    refcount
    dirty / locked / uptodate
    dma_pinned

  Devices[]:
    MMIO registers
    DMA queues
    interrupt vectors
    reset and error state

Kernel
  Task:
    saved context
    kernel stack
    state and wait reason
    address space
    file table

  AddressSpace:
    VMA list
    page table root
    active CPUs that may hold TLB entries

  OpenFile:
    permissions
    file offset or append mode
    reference count
    backing inode, pipe, socket, or device

  Request:
    id/tag
    buffer ownership
    DMA mapping
    status and error code
    wait queue or callback
\end{lstlisting}

用这套模型重新走一次 \code{write(fd, user\_buf, len)}，每一步都能落到对象上：

\begin{longtable}{p{0.15\linewidth}p{0.43\linewidth}p{0.3\linewidth}}
\toprule
动作 & 参与对象 & 不变量 \\
\midrule
进入内核 & CPU、trap frame、task、内核栈 & 用户现场已保存，内核可安全使用自己的栈和寄存器 \\
解析 fd & task、fd table、open file & file 引用有效，权限和偏移语义明确 \\
验证地址 & task、address space、VMA、PTE & 用户范围逐页检查，fault 可修复才继续 \\
接收数据 & page cache 或内核 buffer、physical page & 已接收字节和未接收字节边界明确 \\
提交 I/O & request、device queue、DMA mapping & descriptor、buffer、doorbell 顺序正确 \\
睡眠等待 & task、wait queue、request & task 等的是具体条件，不占 CPU 忙等 \\
完成事件 & device completion、interrupt entry、request & completion 能匹配 request，并记录成功或错误 \\
返回用户态 & task、trap frame、open file、request result & 返回值只反映已经提交到相应语义边界的事实 \\
\bottomrule
\end{longtable}

再把这条路径按时间展开。下面不是新增概念，而是把前面分散讲过的 CPU、MMU、cache、设备和 completion 放到同一条线里：

\begin{longtable}{p{0.12\linewidth}p{0.32\linewidth}p{0.24\linewidth}p{0.2\linewidth}}
\toprule
时间 & 硬件动作 & 内核对象变化 & 读者要抓住什么 \\
\midrule
T0 & 用户态 CPU 执行 \code{syscall} 指令 & trap frame 开始形成，task 仍是当前 task & 入口先保存现场，不先信参数 \\
T1 & CPU 切到内核入口和内核栈 & task 的内核栈承载系统调用路径 & 内核不用用户栈保存状态 \\
T2 & 内核按 fd table 查 open file & file 引用增加，权限和 append 状态被检查 & 小整数变成有生命周期的对象 \\
T3 & CPU 读取用户虚拟地址 & MMU 查 TLB/PTE，可能触发 page fault & 用户指针逐页变成可信字节 \\
T4 & 缺页处理读取 VMA，安装或拒绝 PTE & address space、VMA、page refcount 更新 & not-present 要靠 VMA 解释 \\
T5 & CPU 把字节写入 page cache 或内核 buffer & page 变 dirty、locked 或 uptodate & \code{write} 可能只到内存层 \\
T6 & 回写或 direct I/O 构造设备请求 & request 分配 tag，buffer 建 DMA mapping & 设备工作必须有身份牌 \\
T7 & CPU 写 descriptor，再写 MMIO doorbell & request 进入 submitted 状态 & 顺序必须保证设备先看到描述符 \\
T8 & task 若需等待，离开 CPU & task 挂到 wait queue，调度器运行别人 & 等待不是忙等，也不是丢失调用栈 \\
T9 & 设备按 DMA 地址访问内存并执行命令 & buffer 仍被 request 和 DMA mapping 持有 & completion 前不能释放 buffer \\
T10 & 设备写 completion，触发中断或等待轮询 & driver 按 tag 找 request，记录 status & 中断不是结果，completion 才是结果证据 \\
T11 & 内核唤醒 task，撤销 DMA mapping，释放引用 & task 回到 runnable，request 走向 release & 完成要伴随回收和错误传播 \\
T12 & CPU 恢复 trap frame 返回用户态 & 返回值放入用户可见寄存器 & 用户看到的是接口语义，不是全部硬件细节 \\
\bottomrule
\end{longtable}

这条时间线也能解释为什么 OS 教材里同一个词会在不同章节反复出现。调度看到的是 T8 到 T12：task 为什么能睡下去又醒来。虚拟内存看到的是 T3 到 T4：地址为什么要按页检查。文件系统看到的是 T5 和持久化边界：写入何时对文件可见、何时对崩溃恢复可见。驱动看到的是 T6 到 T11：request 怎样提交，completion 怎样匹配，DMA 怎样收尾。把这些阶段分开后，“硬件重要”就不再是一句口号，而是每个 OS 对象的来源。

再把“提交边界”单独拎出来。OS 里很多 bug 都来自把不同层的完成混成一层：

\begin{longtable}{p{0.2\linewidth}p{0.38\linewidth}p{0.3\linewidth}}
\toprule
完成说法 & 真实含义 & 不能误解成 \\
\midrule
系统调用入口完成 & CPU 已进入内核并保存现场 & 参数已经可信 \\
用户拷贝完成 & 字节已经进入内核 buffer 或 page cache & 数据已经写到设备 \\
request submitted & 描述符和 doorbell 已提交给设备路径 & 设备已经完成 I/O \\
completion 到达 & 设备报告某个 tag 有结果 & 文件系统元数据已经持久 \\
\code{write} 返回 & 对应接口语义下已接收或已提交一定字节 & \code{fsync} 语义已经满足 \\
\code{fsync} 返回 & 文件系统和设备承诺的持久化边界已完成或返回错误 & 所有未来硬件故障都不可能丢数据 \\
\bottomrule
\end{longtable}

所有这些对象还有一个共同生命周期。它们通常不是“创建后用一下就丢”，而是沿着 create、reference、use、wait、complete、release 往前走：

\begin{longtable}{p{0.16\linewidth}p{0.42\linewidth}p{0.3\linewidth}}
\toprule
阶段 & 含义 & 例子 \\
\midrule
create & 把一次硬件事实或 API 请求变成对象 & 建 task、建 VMA、分配 request \\
reference & 防止对象在使用期间被释放 & fd lookup 增加 file 引用，I/O pin 用户页 \\
use & 按对象保存的权限和状态推进工作 & 拷贝用户页、填 descriptor、修改 page cache \\
wait & 条件暂时不满足时让 task 离开 CPU & 等页读入、等锁、等 completion \\
complete & 外部事实回来后记录结果并唤醒等待者 & page fault resolved、DMA done、journal committed \\
release & 最后一个引用消失后撤销映射和释放资源 & put file、unpin page、free request \\
\bottomrule
\end{longtable}

这个生命周期能防住很多直觉错误。fd 查找之后要给 file 增加引用，是因为另一个线程可能马上 \code{close(fd)}；直接 I/O 要 pin page，是因为用户可能马上 \code{munmap} 或退出；request 完成后不能只唤醒 task，还要撤销 DMA mapping，是因为设备地址空间也曾经持有访问权；VMA 删除后要处理 TLB，是因为 CPU 可能还握着旧翻译。OS 对象的引用计数、锁、等待队列和状态位，都是在给这些生命周期边界补证据。

用一个排障案例把全章连起来。用户报告：

\begin{lstlisting}
write(fd, buf, len) sometimes hangs
\end{lstlisting}

不能只盯着 \code{write} 这一行。按本章模型，要顺着对象链往下查：

\begin{longtable}{p{0.18\linewidth}p{0.39\linewidth}p{0.31\linewidth}}
\toprule
检查点 & 可能发现 & 指向的章节 \\
\midrule
task 状态 & task 是 running、runnable，还是 blocked & 调度与等待 \\
等待原因 & task 挂在哪个 wait queue 上 & 锁、I/O、内存回收或页 fault \\
用户拷贝 & 是否卡在 page fault、缺页读入或内存压力 & 虚拟内存 \\
file 对象 & fd 是否指向慢设备、管道、socket 或普通文件 & 文件与 VFS \\
page cache & 页是否 locked、writeback、dirty 堆积 & 文件缓存与回写 \\
request & 块请求是否 submitted、timeout 或等待队列深 & 块层和驱动 \\
completion & 设备是否写了 completion，中断是否送达或被轮询收割 & 中断与设备 \\
DMA mapping & buffer 是否还被设备持有，IOMMU 是否报错 & DMA 与硬件隔离 \\
\bottomrule
\end{longtable}

假设排查发现 task 睡在某个 request 的 wait queue 上。下一步不是说“磁盘慢”就结束，而是继续看 request state。如果 request 还没 submitted，问题可能在块层限流、队列冻结或内存分配；如果已经 submitted 但没有 completion，问题可能在设备队列、doorbell、DMA 地址、设备固件或中断；如果 completion 已到但 task 没醒，问题可能在驱动收割、request tag 匹配或 wait queue 唤醒；如果 task 醒了但 \code{write} 仍不返回，问题可能在文件系统日志提交或错误回报。每一步都对应一个对象，不靠猜。

再换一个现象：

\begin{lstlisting}
write(fd, bad_user_ptr, 4096) returns EFAULT
\end{lstlisting}

这里设备路径可能完全没有发生。CPU 在用户拷贝阶段访问 \code{bad\_user\_ptr}，page fault 入口记录 fault address，内核查当前 task 的 address space，发现没有 VMA 或权限不对，于是系统调用返回 \code{EFAULT}。如果程序传的地址第一半合法、第二半非法，可能已经复制一部分；不同对象类型对“部分成功后遇到错误”有不同语义。这条路径让人看到：同样叫 \code{write}，有些失败发生在 CPU/MMU 边界，甚至还没碰到文件系统和设备。

再看性能案例：

\begin{lstlisting}
small writes are fast, fsync is slow
\end{lstlisting}

这通常不是矛盾。小 \code{write} 可能只把字节放入 page cache 并标 dirty，返回时没有等待设备持久化；\code{fsync} 要把此前积累的脏页、inode 元数据和日志事务推进到持久化边界，还可能等待设备 flush。若设备有易失缓存，completion 和持久化之间还隔着 flush/FUA 语义。这里 page cache、writeback request、journal transaction 和 device completion 都是不同对象，不能用一个“写完了”概括。

这张表会贯穿后文。调度要区分“task 可运行”和“task 正在 CPU 上运行”；内存要区分“VMA 合法”和“PTE 已 present”；文件要区分“page cache dirty”和“设备持久化”；网络要区分“包进入网卡 DMA buffer”和“用户 \code{recv} 已拿到字节”；安全要区分“对象存在”和“当前主体有权使用它”。

把第三章压成后文索引，可以得到下面这张映射表：

\begin{longtable}{p{0.2\linewidth}p{0.32\linewidth}p{0.32\linewidth}}
\toprule
硬件事实 & 内核对象 & 后文会展开的问题 \\
\midrule
CPU 被同步或异步入口打断 & trap frame、task、kernel stack & 系统调用、异常、中断、调度 \\
虚拟地址要翻译和检查权限 & address space、VMA、PTE、TLB state & 缺页、COW、mmap、用户拷贝 \\
物理页会被共享、锁住、写脏或 DMA 引用 & page、page cache、pin/refcount & 内存分配、回收、回写、直接 I/O \\
fd 只是进程表里的索引 & fd table、open file、inode/socket/pipe & VFS、权限、偏移、并发 close \\
设备通过队列和寄存器协作 & driver、queue、descriptor、doorbell & 块层、网络、字符设备、驱动模型 \\
完成异步到达且可能失败 & request、completion、wait queue、timeout & 睡眠唤醒、超时、取消、reset \\
持久化不是普通内存写 & dirty page、journal transaction、flush request & 文件系统一致性、fsync、崩溃恢复 \\
\bottomrule
\end{longtable}

最后给一份阅读 OS 机制的检查表。以后遇到任何机制，都按这六问拆：

\begin{longtable}{p{0.18\linewidth}p{0.42\linewidth}p{0.28\linewidth}}
\toprule
问题 & 要找的证据 & 例子 \\
\midrule
状态在哪里 & 哪个对象保存长期事实，不在 CPU 瞬时寄存器里 & task、VMA、request、inode \\
入口在哪里 & 硬件或 API 从哪个受控点进入内核 & syscall、page fault、IRQ、timer \\
权限在哪里检查 & 谁把不可信参数变成可用对象 & fd lookup、VMA check、LSM、IOMMU \\
等待在哪里表达 & task 睡在哪个条件上，谁负责唤醒 & wait queue、completion、futex queue \\
提交点在哪里 & 哪一步之后对上层语义可见 & PTE install、file size update、doorbell、journal commit \\
失败怎样收敛 & 哪些对象回滚、释放、保留错误码或重试 & errno、request error、page refcount、reset recovery \\
\bottomrule
\end{longtable}

\begin{keyidea}
硬件不给应用直接共享整台机器的安全办法。CPU、MMU、cache、设备和持久化介质只暴露低层状态；内核必须把这些状态组织成 task、address space、open file、request、wait queue、completion 和恢复证据。后面每一章都在细化这套转换。
\end{keyidea}
